\documentclass[11pt]{article}

\usepackage[margin=1in]{geometry}
\usepackage{amsmath,amssymb}
\usepackage{booktabs}
\usepackage{siunitx}

\title{Extended Vehicular A (EVA) Channel Model}
\author{}
\date{}

\begin{document}
\maketitle

\section{Overview}
The Extended Vehicular A (EVA) channel is one of the standard multipath
fading channel models specified by 3GPP in TS~36.101 for the conformance
testing of LTE and LTE-Advanced user equipment. It represents a typical
vehicular propagation environment and is widely used in the simulation of
mobile broadband systems.

The channel is modeled as a tapped delay line (TDL) with $L = 9$ resolvable
paths. Each path is characterized by a complex path gain $\alpha_\ell$, an
excess delay $\tau_\ell$, and its own Doppler shift $\nu_\ell$. The
time-varying baseband channel impulse response is
\begin{equation}
    h(t,\tau) = \sum_{\ell=0}^{L-1} \alpha_\ell\,e^{j2\pi\nu_\ell t}\,\delta(\tau - \tau_\ell),
\end{equation}
where $\alpha_\ell$ is a zero-mean complex Gaussian random variable whose
variance equals the relative average power $P_\ell$ of the $\ell$-th tap,
and $\nu_\ell$ is bounded by the maximum Doppler frequency,
$|\nu_\ell| \le f_d$.

For a transmitted signal $x(t)$, the received signal $y(t)$ is obtained as
the convolution of $x(t)$ with $h(t,\tau)$, plus additive white Gaussian
noise $n(t)$:
\begin{equation}
    y(t) = \int_{-\infty}^{\infty} h(t,\tau)\,x(t-\tau)\,\mathrm{d}\tau + n(t)
         = \sum_{\ell=0}^{L-1} \alpha_\ell\,e^{j2\pi\nu_\ell t}\,x(t-\tau_\ell) + n(t).
\end{equation}

\section{Power Delay Profile}
The excess tap delays and relative average powers of the EVA channel are
given in Table~\ref{tab:eva}.

\begin{table}[h]
    \centering
    \caption{EVA channel power delay profile (3GPP TS~36.101).}
    \label{tab:eva}
    \begin{tabular}{c S[table-format=4.0] S[table-format=-2.1]}
        \toprule
        Tap index $\ell$ & {Excess delay $\tau_\ell$ (ns)} & {Relative power $P_\ell$ (dB)} \\
        \midrule
        0 &    0 &   0.0 \\
        1 &   30 &  -1.5 \\
        2 &  150 &  -1.4 \\
        3 &  310 &  -3.6 \\
        4 &  370 &  -0.6 \\
        5 &  710 &  -9.1 \\
        6 & 1090 &  -7.0 \\
        7 & 1730 & -12.0 \\
        8 & 2510 & -16.9 \\
        \bottomrule
    \end{tabular}
\end{table}

The maximum excess delay is $\tau_{\max} = \SI{2.51}{\micro\second}$, and the
RMS delay spread is approximately $\SI{357}{\nano\second}$.

\section{Doppler Configurations}
Following the classical (Jakes) scattering assumption, the Doppler shift of
each path is determined by the angle of arrival $\theta_\ell$ between the
direction of motion and the incoming wavefront:
\begin{equation}
    \nu_\ell = \nu_{\max}\cos(\theta_\ell),
    \qquad \theta_\ell \sim \mathcal{U}[0,\,2\pi),
\end{equation}
where $\nu_{\max} = v\,f_c / c$ is the maximum Doppler frequency, $v$ is the
UE speed, $f_c$ is the carrier frequency, and $c$ is the speed of light.
The angles $\{\theta_\ell\}$ are drawn independently across paths, so each
tap experiences a different Doppler shift in $[-\nu_{\max},\,\nu_{\max}]$.

Assuming a UE speed of $v = \SI{120}{\kilo\meter\per\hour} \approx
\SI{33.33}{\meter\per\second}$ and an LTE carrier frequency of
$f_c = \SI{2}{\giga\hertz}$, the maximum Doppler frequency is
\begin{equation}
    \nu_{\max}
    = \frac{v\,f_c}{c}
    = \frac{33.33 \times 2 \times 10^{9}}{3 \times 10^{8}}
    \approx \SI{222.2}{\hertz}.
\end{equation}

\section{Reduced EVA Channel}
For low-complexity simulation and analysis, a reduced version of the EVA
channel is often used in which only the three most powerful taps of the
original profile are retained. From Table~\ref{tab:eva}, these correspond to
path indices $\ell \in \{0,\,4,\,2\}$ (listed in order of decreasing
average power), giving the reduced power delay profile in
Table~\ref{tab:eva-reduced}.

\begin{table}[h]
    \centering
    \caption{Reduced EVA channel power delay profile (three strongest taps).}
    \label{tab:eva-reduced}
    \begin{tabular}{c c S[table-format=3.0] S[table-format=-1.1]}
        \toprule
        Tap index $\ell$ & Original index & {Excess delay $\tau_\ell$ (ns)} & {Relative power $P_\ell$ (dB)} \\
        \midrule
        0 & 0 &   0 &  0.0 \\
        1 & 4 & 370 & -0.6 \\
        2 & 2 & 150 & -1.4 \\
        \bottomrule
    \end{tabular}
\end{table}

With $L = 3$, the channel impulse response and received signal retain the
same form as before:
\begin{align}
    h(t,\tau) &= \sum_{\ell=0}^{2} \alpha_\ell\,e^{j2\pi\nu_\ell t}\,\delta(\tau - \tau_\ell), \\
    y(t)      &= \sum_{\ell=0}^{2} \alpha_\ell\,e^{j2\pi\nu_\ell t}\,x(t-\tau_\ell) + n(t),
\end{align}
where the path parameters $\{\alpha_\ell, \tau_\ell, \nu_\ell\}$ are taken
from the reduced profile. To preserve unit average channel energy, the tap
powers may optionally be renormalized so that
$\sum_{\ell=0}^{2}\mathbb{E}\{|\alpha_\ell|^2\} = 1$.

\end{document}
